\subsection{Ideals in a ring of fractions}

PROPOSITION 11. Let A be a ring and S a multiplicative subset of A. For every ideal \(\mathfrak{b}'\) of \(\mathbf{S}^{-1}\mathbf{A}\) , let \(\mathfrak{b} = (i_{\mathbf{A}}^{\mathbf{S}})^{-1}(\mathfrak{b}')\) be such that \(\mathfrak{b}' = \mathbf{S}^{-1}\mathfrak{b}\) .

\begin{enumerate}
\def\labelenumi{(\roman{enumi})}
\item
  Let \(f\) be the canonical homomorphism \(\mathbf{A} \to \mathbf{A} / \mathfrak{b}\) . The homomorphism from \(\mathbf{S}^{-1}\mathbf{A}\) to \((f(\mathbf{S}))^{-1}(\mathbf{A} / \mathfrak{b})\) canonically associated with \(f\) (no. 1, Proposition 2) is surjective and its kernel is \(\mathfrak{b}'\) , which defines, by taking quotients, a canonical isomorphism of \((\mathbf{S}^{-1}\mathbf{A}) / \mathfrak{b}'\) onto \((f(\mathbf{S}))^{-1}(\mathbf{A} / \mathfrak{b})\) . Moreover, the canonical homomorphism from \(\mathbf{A} / \mathfrak{b}\) to \((f(\mathbf{S}))^{-1}(\mathbf{A} / \mathfrak{b})\) is injective.
\item
  The mapping \(\mathfrak{b}' \mapsto \mathfrak{b} = (i_{\mathbf{M}}^{\mathbf{S}})^{-1}(\mathfrak{b}')\) , restricted to the set of maximal (resp. prime) ideals of \(\mathbf{S}^{-1}\mathbf{A}\) , is an isomorphism (with respect to inclusion) of this set onto the set of ideals of \(\mathbf{A}\) which are maximal among those which do not meet \(\mathbf{S}\) (resp. the set of prime ideals of \(\mathbf{A}\) not meeting \(\mathbf{S}\) ).
\item
  If \(\mathfrak{q}'\) is a prime ideal of \(S^{-1}A\) and \(\mathfrak{q} = (i_A^S)^{-1}(\mathfrak{q}')\) , there exists an isomorphism of the ring of fractions \(A_{\mathfrak{q}}\) onto the ring \((S^{-1}A)_{\mathfrak{q}}\) , which maps \(a/b\) to \((a/1)/(b/1)\) , where \(a \in A\) , \(b \in A - \mathfrak{q}\) .

\end{enumerate}

\begin{enumerate}
\def\labelenumi{(\roman{enumi})}
\item
  \((f(S))^{-1}(A/b)\) can be identified with \(S^{-1}(A/b)\) by means of the canonical isomorphism between these two modules (no. 2, Proposition 6). The exact sequence \(0 \rightarrow b \rightarrow A \rightarrow A/b \rightarrow 0\) then induces an exact sequence

\[
0 \rightarrow \mathrm{S} ^ {- 1} \mathfrak {b} \rightarrow \mathrm{S} ^ {- 1} \mathrm{A} \rightarrow \mathrm{S} ^ {- 1} (\mathrm{A} / \mathfrak {b}) \rightarrow 0
\]

(no. 4, Theorem 1) whose existence proves the first assertion of (i), taking account of the fact that \(\mathfrak{b}' = \mathrm{S}^{-1}\mathfrak{b}\) . Since \(\mathfrak{b}\) is saturated with respect S, the conditions \(a \in \mathbf{A}\) , \(s \in \mathbf{S}\) , \(as \in \mathfrak{b}\) imply \(a \in \mathfrak{b}\) ; the homothety of ratio \(s\) on \(\mathbf{A}/\mathfrak{b}\) is then injective, which proves the second assertion of (i).

\item
  We note first that the relation \(b' = S^{-1}A\) is equivalent to the relation \(b \cap S \neq \varnothing\) , the latter expressing the fact that \(b'\) contains invertible elements of \(S^{-1}A\) . It follows from no. 4, Proposition 10 (iii) that \(b' \mapsto b = (i_{\mathrm{A}}^{\mathrm{S}})^{-1}(b')\) is an isomorphism (with respect to inclusion) of the set of ideals of \(S^{-1}A\) distinct from \(S^{-1}A\) onto the set F of ideals of A not meeting S and satisfying condition (MS) of Proposition 10. If \(b'\) is maximal (resp. prime), clearly \(b'\) is maximal in F (resp. prime) and conversely (by (i)). On the other hand, if r is an ideal of A disjoint from S, its saturation \(r_{1}\) with respect to S is an ideal of A containing r and disjoint from S: for no element \(a \in S\) can satisfy \(sa \in r\) for some \(s \in S\) , since it would follow that \(sa \in r \cap S\) . We conclude that, if r is maximal among the ideals of A meeting S, it is maximal in F. Similarly, if r is a prime ideal not meeting S, it satisfies condition (MS) of no. 4, Proposition 10 by definition of prime ideals and hence belongs to F This completes the proof of (ii).
\item
  Suppose that \(q'\) is prime and such that q is also prime. The set \(T = A - q\) is a multiplicative subset of A which contains S, whence \(ST = T\) . We write \(T' = i_{\mathbf{A}}^{\mathbf{S}}(\mathbf{T})\) ; it follows from no. 3, Proposition 7 (i) that there exists a unique isomorphism j of \(T^{-1}A = A_{q}\) onto \(T'^{-1}(S^{-1}A)\) such that

\[
j (a / b) = (a / 1) / (b / 1),
\]

where \(a \in \mathbf{A}\) and \(b \in \mathbf{T}\) . On the other hand \(\mathbf{T}'\) obviously does not meet \(\mathfrak{q}'\) ; conversely, let \(a / s \in \mathrm{S}^{-1}\mathrm{A}\) ; since \(1 / s\) is invertible in \(\mathrm{S}^{-1}\mathrm{A}\) , the condition \(a / s \notin \mathfrak{q}'\) is equivalent to \(i_{\mathbf{A}}^{\mathbf{S}}(a) = a / 1 \notin \mathfrak{q}'\) and hence to \(a \notin \mathfrak{q}\) ; it follows that \(\mathrm{S}^{-1}\mathrm{A} - \mathfrak{q}' = \mathrm{S}^{-1}\mathrm{T}'\) and hence, by Proposition 8 of no. 3, \(\mathrm{T}'^{-1}(\mathrm{S}^{-1}\mathrm{A}) = (\mathrm{S}^{-1}\mathrm{A})_{\mathfrak{q}'}\) .

The isomorphism defined in (iii) is called canonical. If A is an integral domain, the canonical isomorphisms of \(A_{q}\) and \((\mathrm{S}^{-1}\mathrm{A})_{q'}\) onto subrings of the field of fractions K of A have the same image.
\end{enumerate}

Remark. For an ideal a of A to satisfy \(S^{-1}a = S^{-1}A\) (or, what amounts to the same thing by no. 4, Theorem 1, \(S^{-1}(A/a) = 0\) ), it is necessary and sufficient that \(a \cap S \neq \varnothing\) , as follows immediately from the definitions.

COROLLARY 1. Let A be a ring and S a multiplicative subset of A. Every ideal p of A which is maximal among those which do not meet S is prime.

By Proposition 11, the hypothesis on p means that \(\mathfrak{p} = (i_{\mathrm{A}}^{\mathrm{S}})^{-1}(\mathfrak{m}')\) , where \(m'\) is a maximal ideal of \(S^{-1}A\) ; as \(m'\) is prime, so is p.

COROLLARY 2. Let A be a ring and S a multiplicative subset of A. For every ideal a of A not meeting S there exists a prime ideal containing a and not meeting S.

\(S^{-1}a \neq S^{-1}A\) (Remark) and hence there exists a maximal ideal of \(S^{-1}A\) containing \(S^{-1}a\) (Algebra, Chapter I, § 8, no. 7, Theorem 2) and the corollary follows from Proposition 11 (ii).

COROLLARY 3. Let A, B be two rings, \(\rho\) a homomorphism from A to B and \(\mathfrak{p}\) a prime ideal of A. For there to exist a prime ideal \(\mathfrak{p}'\) of B such that \(\overline{\rho}^{1}(\mathfrak{p}') = \mathfrak{p}\) , it is necessary and sufficient that \(\overline{\rho}^{1}(\mathrm{B}\rho(\mathfrak{p})) = \mathfrak{p}\) .

If there exists an ideal \(p'\) of \(B\) such that \(\overline{\rho}^{1}(p') = p\) , then \(\rho(p) \subset p'\) , whence \(B\rho(p) \subset p'\) and \(\overline{\rho}^{1}(B\rho(p)) \subset \overline{\rho}^{1}(p') \subset p\) ; as the converse inclusion is obvious, \(\overline{\sigma}^{1}(B\rho(p)) = p\) . Conversely, suppose that \(\overline{\rho}^{1}(B\rho(p)) = p\) and consider the multiplicative subset \(S = \rho(A - p)\) of \(B\) ; the hypothesis shows that

\[
\mathrm{S} \cap \mathrm{B} _ {\rho} (\mathfrak {p}) = \varnothing ;
\]

by Corollary 2 there exists a prime ideal \(\mathfrak{p}'\) of B containing \(\mathrm{B}\varphi(\mathfrak{p})\) and not meeting S; then \(\overline{\rho}^{1}(\mathfrak{p}')\) contains \(\mathfrak{p}\) and cannot contain any element of A - \(\mathfrak{p}\) and hence is equal to \(\mathfrak{p}\) .

COROLLARY 4. Let A, B be two rings and \(\rho\) a homomorphism from A to B.

\begin{enumerate}
\def\labelenumi{(\roman{enumi})}
\item
  Suppose that there exists a B-module E such that \(\rho_{*}(\mathbf{E})\) is a faithfully flat A-module. Then, for every prime ideal \(\mathfrak{p}\) of A, there exists a prime ideal \(\mathfrak{p}'\) of B such that \(\overline{\rho}^{-1}(\mathfrak{p}') = \mathfrak{p}\) .
\item
  Conversely, suppose that B is a flat A-module. Then, if, for every prime ideal p of A, there exists an ideal \(p'\) of B such that \(\overline{\rho}^{-1}(p') = p\) , B is a faithfully flat A-module.
\end{enumerate}

\begin{enumerate}
\def\labelenumi{(\roman{enumi})}
\item
  The hypothesis implies that, for every ideal \(\mathfrak{a}\) of \(\mathbf{A}\) , \(\bar{\rho}^{1}(\mathrm{B}\varphi (\mathfrak{a})) = \mathfrak{a}\) (Chapter I, § 3, no. 5, Proposition 8 (ii)) and it is sufficient to apply Corollary 3.

\item
  It is sufficient to show that, for every maximal ideal m of A, there exists a maximal ideal \(m'\) of B such that \(\bar{\rho}^{-1}(m') = m\) (Chapter I, §2, no. 5, Proposition 9 (e)). Now there exists by hypothesis an ideal q of B such that \(\bar{\rho}^{-1}(q) = m\) ; as \(q \neq B\) , there exists a maximal ideal \(m'\) of B containing q and consequently \(\bar{\rho}^{-1}(m') \supset m\) ; but, as \(\bar{\rho}^{-1}(m')\) cannot contain 1, \(\bar{\rho}^{-1}(m') = m\) .
\end{enumerate}

COROLLARY 5. Let A be a ring, S a multiplicative subset of A and B a ring such that \(i_{\mathbf{A}}^{\mathbf{S}}(\mathbf{A}) \subset \mathbf{B} \subset \mathbf{S}^{-1}\mathbf{A}\) . Let \(\mathfrak{q}\) be a prime ideal of B such that the prime ideal \(\mathfrak{p} = (i_{\mathbf{A}}^{\mathbf{S}})^{-1}(\mathfrak{q})\) of A does not meet S and let \(\mathfrak{p}'\) be the prime ideal \(\mathbf{S}^{-1}\mathfrak{p}\) of \(\mathbf{S}^{-1}\mathbf{A}\) . Then \(\mathfrak{p}' \cap \mathbf{B} = \mathfrak{q}\) .

Let \(S' = i_{\mathbf{A}}^{\mathbf{S}}(\mathbf{S})\) ; a canonical isomorphism has been defined from \(S'^{-1}\mathbf{B}\) to \(S^{-1}\mathbf{A}\) (no. 1, Corollary 4 to Proposition 2); we identify these two rings by means of this isomorphism. As \(q \cap S' = \varnothing\) , \(q' = S'^{-1}q\) is the unique prime ideal of \(S^{-1}\mathbf{A} = S'^{-1}\mathbf{B}\) such that \(q' \cap B = (i_{\mathbf{A}}^{\mathbf{S}})^{-1}(q') = q\) (Proposition 11 (ii)), whence \((i_{\mathbf{A}}^{\mathbf{S}})^{-1}(q') = p\) ; consequently \(q' = p'\) (Proposition 11 (ii)).

In the notation of Corollary 5, there are canonical isomorphisms of \(\mathbf{A}_{\mathfrak{p}}\) and \(\mathbf{B}_{\mathfrak{q}}\) onto \((\mathrm{S}^{-1}\mathrm{A})_{\mathfrak{p}}\) . (Proposition 11 (iii)) and hence a canonical isomorphism \(\mathbf{A}_{\mathfrak{p}} \to \mathbf{B}_{\mathfrak{q}}\) .

